\documentclass[12pt]{article}
\usepackage[utf8]{inputenc}
\usepackage[T1]{fontenc}
\usepackage{hyperref}
\title{Tumbuka Language Resources for K-12 Curriculum Development}
\author{Compiled by Claude Sonnet 5 (claude-sonnet-5) for Andr\'as Kornai}
\date{}
\begin{document}
\maketitle

\section{Introduction}
Tumbuka (Chitumbuka, Citumbuka; ISO 639-3 tum) is spoken in northern Malawi and eastern Zambia. The search returned mostly bibliographic references and a few online pages. Items are listed as found. Availability and current editions were not verified, and many sources have no URL.

\section{Primary Educational Resources}
The search found no school textbooks or primers with URLs. It did turn up two learner and teaching-oriented items, listed below.
\begin{itemize}
\item Duncan McNicholl, \emph{The No-Nonsense Guide to Learning Chitumbuka: Volume 1} (2010), a learner's guide. It was mentioned in the search summary only, with no URL.
\item Free Tumbuka Phrasebook. \url{https://freelanguage.org/node/77453}
\item Tumbuka phrasebook (Wikivoyage), a traveller's phrasebook. \url{https://en.wikivoyage.org/wiki/Tumbuka_phrasebook}
\end{itemize}

\section{Linguistic and Reference Materials}
\begin{itemize}
\item \emph{Mung'anamulira mazgo wa Citumbuka}, a Citumbuka dictionary (2018), Centre for Language Studies, University of Malawi. Pascal Kishindo is named as editor in the search summary. No URL was found.
\item W.\,M. Turner, \emph{Tumbuka--Tonga--English Dictionary} (1952, Hetherwick Press, Blantyre). No URL was found.
\item Walter Angus Elmslie, \emph{Introductory Grammar of the Tumbuka Language} (1923, Livingstonia Mission Press). Glottolog reference: \url{https://glottolog.org/resource/reference/id/57905}
\item Kishindo and Lipenga, \emph{Parlons citumbuka} (L'Harmattan, 2006), a French-language introduction. No URL was found.
\item Wikipedia, Tumbuka grammar. \url{https://en.wikipedia.org/wiki/Tumbuka_grammar}
\item Wikipedia, Tumbuka orthography. \url{https://en.wikipedia.org/wiki/Tumbuka_orthography}
\item WEBBALL bibliographic entry. \url{https://reflex.cnrs.fr/Lexiques/webball/biblio_bas.php?id=20221}
\end{itemize}

\section{Religious and Community Resources}
No scripture or community resources appeared in the results.

\section{Digital and Audio Resources}
No audio, video or software resources were found in this search, apart from the online phrasebooks listed above.

\section{Curriculum Development Context}
Tumbuka has a documented orthography, historical grammars and a recent dictionary from the University of Malawi, so there is a base for curriculum work. The search did not confirm whether Tumbuka is currently taught in schools, or which materials the Malawi and Zambia education systems use. Those points need follow-up with the Centre for Language Studies and the ministries of education.

\end{document}